\scanpage{120}
\uline{Sur une certaine fonction holomorphe \deleted{\unsure{de variable|définie}} \inserted{sur l'}espace de Banach $l^1$.}

Soit, pour $\lambda=(\lambda_i)\in l^1$ et tout entier $n>0$
\[
\alpha_n(\lambda)=\sum_{i_1<\mathord{-}\mathord{.}\mathord{.}<i_n}\lambda_{i_1}\mathord{.}\mathord{-}\lambda_{i_n}
\]
On a
\[
\alpha_n(\lambda)\leq\frac{\|\lambda\|_1^n}{n!}
\]
les $\alpha_n(\lambda)$ sont les coefficients de la fonction entière
de genre zéro
\[
f(\unsure{n|\lambda})=\prod_1(1+\lambda_i z)
\]
et on obtient, à un facteur \unsure{constant} près, les
coefficients les plus généraux d'une fonction entière
de genre zéro.

\uline{Proposition} Pour tout $\lambda\in l^1$, la série $\sum_0^\infty n!\alpha_n(\lambda)$
\margin{\tikz[x=1pt,y=1pt]{\draw[line width=.45pt] (0,0)--(.6,-32)--(2,-67)--(3.2,-101);}}
est absolument convergente, soit $S(\lambda)$ sa somme.
La fonction $\lambda\to S(\lambda)$ est holomorphe au voisinage
de $0$, avec un rayon de convergence égal à $1$,
\deleted{\unsure{ce qui fait que ses|ce qui fait que sur}} \inserted{et ses restrictions} aux sous-espaces vectoriels
complexes de dim finie sont holomorphes.
Cependant, \deleted{elle} cette fonction n'est pas bornée
dans la boule unité de $l^1$, et n'est par
suite pas une fonction entière de $\lambda$.

\uline{Corollaire} On a \inserted{pour $\lambda\in l^1$} $\sqrt[n]{n!|\alpha_n(\lambda)|}\to 0$,
\margin{\tikz[x=1pt,y=1pt]{\draw[line width=.45pt] (0,0)--(-.5,-14)--(.4,-38)--(1.5,-69);}}
$\sqrt[n]{|\alpha_n(\lambda)|}=o\left(\frac{1}{n}\right)$, $\alpha_n(\lambda)=\frac{\varepsilon_n^n}{n!}$ avec $\varepsilon_n\to 0$
Si $\sum a_nz^n$ est une fonction entière de
genre zéro, la suite des $a_n$ jouit
des $=\mathrm s$ propriétés

En effet, $\sqrt[n]{n!|\alpha_n(\lambda)|}\to 0$ exprime que
\margin{121}
